% !TEX root = ../ComputationalOTFnT.tex

\chapter{変分ワッサーシュタイン問題}
\label{c-variational} % a label for the chapter, to refer to it later

データ解析において、確率測度間の一般的なダイバージェンス（\emph{例：}ユークリッド、全変動、ヘリンガー、カルバック・ライブラー）は、パラメータ推定問題におけるフィッティング誤差や損失の測定によく用いられる。本章に至るまで、我々は最適輸送幾何学が他の情報ダイバージェンスにはない固有の能力、すなわち物理的な概念（質量の移動）と幾何学（観測値やビン間のグラウンドコスト）を活用して測度を比較する能力を持つことを論じてきた。これら二つの事実を組み合わせると、ワッサーシュタイン距離を損失関数として使用することは非常に魅力的となる。このアイデアは近年、様々な応用問題に対して探究されてきた。しかし、このアイデアに伴う主な技術的課題は、ワッサーシュタイン距離を効率的に近似し微分することにある。本章ではいくつかの動機付けとなる例から始め、本書の前半の章で提示した様々な数値スキームが変分ワッサーシュタイン問題を解くためにどのように利用できるかを示す。

画像処理において、ワッサーシュタイン距離はテクスチャ合成の損失として使用でき~\citep{2016-tartavel-siims}、合成されたサンプルと入力サンプルの統計量間の不一致を考慮する。また、画像領域の統計的均質性を考慮した画像セグメンテーションにも使用される~\citep{SchnorSegmentation,RabinPapadakisSSVM,peyre2012wasserstein,ni2009local,schmitzer2013object,liimage}。
%
ワッサーシュタイン距離はまた、測定値が確率測度である逆問題に対する非常に自然な適合度項でもある。例えば、画像復元~\citep{lellmann2014imaging}、トモグラフィック反転~\citep{AbrahamRadon}、密度正則化~\citep{Burger-JKO}、粒子画像流速計測~\citep{saumier2015optimal}、スパース回復と圧縮センシング~\citep{indyk2011k}、および地震波インバージョン~\citep{metivier2016optimal}がある。
%
測度間の距離（主に\S\ref{sec-mmd}で示されるカーネルベースのもの）は、計算解剖学~\citep{vaillant2005surface}における形状マッチング（リフトされた空間上の測度として表現され、しばしばカレントと呼ばれる）に日常的に使用されるが、最適輸送距離は興味深い代替手段を提供する~\citep{2017-feydy-miccai}。
%
ヒストグラムのデータセットの次元を削減するために、\citeauthor{lee1999learning}は非負行列分解問題が再構成損失を定量化するためにカルバック・ライブラーダイバージェンスを用いて定式化できることを示した~\citep{lee1999learning}。これらのヒストグラムのビンの幾何学に関する事前情報が利用可能な場合、ワッサーシュタイン距離を代わりに使用でき、著しく異なる結果が得られる~\citep{Sandler09,ZenICPR14,pmlr-v51-rolet16}。

これらの問題はすべて、ワッサーシュタイン距離の勾配、またはその近似へのアクセスを必要とするという共通点がある。本節ではまずそのような勾配を近似する方法を提示し、次に変分ワッサーシュタイン問題として定式化できる三つの重要な応用を紹介する。

%%%%%%%%%%%%%%%%%%%%%%%%%%%%%%%%%%%%%
\section{ワッサーシュタイン損失の微分}

統計学、テキスト処理、または画像処理において、通常、測定から得られる確率分布$\be$を、パラメータ化された分布の族$\{\al_\th,\th\in\Theta\}$と比較する必要がある。ここで$\Theta$はユークリッド空間の部分集合である。このような比較は「損失」または「適合度」項を通じて行われ、本節ではそれがワッサーシュタイン距離である。
%
最も単純なシナリオでは、適切なパラメータ$\th$の計算は直接的な最小化により得られる
\eql{\label{eq-wloss-generic}
	\umin{\th \in \Theta} \Ee(\th) \eqdef \MK_\c(\al_\th,\be).
}
もちろん、より複雑な問題も考えられる。例えば、\S\ref{sec-bary}で述べる重心問題はこのような項の和からなる。しかし、これらのより高度な問題のほとんどは、連鎖律を使用して陽に微分を計算するか、\S\ref{rem-auto-diff}で提唱される自動微分を使用することで、上記の基本的な場合に定義されたツールを適用して解くことができる。

\paragraph{凸性。} 二つのヒストグラムまたは二つの密度間のワッサーシュタイン距離は、~\eqref{eq-dual}および~\eqref{eq-dual-generic}で示されるように、その二つの入力に関して凸である。したがって、パラメータ$\theta$自体がヒストグラムである場合、すなわち$\Theta=\simplex_n$かつ$\al_\theta=\theta$である場合、またはより一般に$\theta$がシンプレックス内の$K$個の重みを記述する場合、$\Theta=\simplex_K$とし$\alpha_\theta=\sum_{i=1}^K \theta_i \alpha_i$が$\simplex_N$内の既知のアトム$\alpha_1,\dots,\alpha_K$の凸結合である場合、問題~\eqref{eq-wloss-generic}は凸のままである（前者は重心問題に対応し、後者はワッサーシュタイン損失を用いた辞書学習問題の一反復に対応する~\citep{pmlr-v51-rolet16}）。しかし、より一般的なパラメータ化$\th\mapsto \al_\th$に対しては、問題~\eqref{eq-wloss-generic}は一般に凸ではない。


\paragraph{単純な場合。} ワッサーシュタイン距離が閉じた形で得られる単純な場合、例えば一変量（\S\ref{rem-1d-ot-generic}参照）または楕円等高分布（\S\ref{rem-dist-gaussians}参照）の場合には、簡単な回避策が存在する。それらは主に、ワッサーシュタイン距離をこれらの分布の適切な表現間のより単純な距離（一変量測度に対する分位関数上のユークリッド距離、同じ族の楕円等高分布に対する共分散行列上のビュール計量）として表現し、問題~\eqref{eq-wloss-generic}をそのような表現上で直接解くことからなる。

しかしほとんどの場合、問題~\eqref{eq-wloss-generic}の局所最小解を計算するためには、$\al_\th$の慎重な離散化に頼る必要がある。二つのアプローチが考えられる：オイラー的またはラグランジュ的である。図~\ref{fig-eulerian-lagrangian}はこれら二つの基本的な離散化スキームの違いを示している。議論を過度に単純化する危険を冒して言えば、オイラー離散化は測度が低次元空間上に台を持つ場合（形状や色空間を扱う場合など）、または本質的に離散的な問題（文字列やテキスト解析から生じるものなど）に最も適している。観測値が高次元空間で連続値をとるフィッティング問題に適用する場合、ラグランジュ的な視点が通常唯一の適切な選択である。
%were already discussed in~\S\ref{sec-grad-flows} above in the context of gradient flows discretization, and we use the same strategies here.
%


\begin{figure}[h!]
\centering
\includegraphics[width=\linewidth]{eulerian-lagrangian/eulerian-lagrangian.pdf}
\caption{\label{fig-eulerian-lagrangian}
密度を持つ連続分布（紫、左）の、ラグランジュ表現$\frac{1}{n}\sum_i \de_{x_i}$（青、上）と、サイズが増加するグリッド上のセルを表す$x_i$を用いたオイラー表現$\sum_i \a_i \de_{x_i}$（赤、下）による次第に細かくなる離散化。オイラー的な視点ではピクセル化された画像から始まり、解像度が十分に細かくなるとほぼ元の密度と一致するようになる。重み$\a_i$は各ピクセルセルの強度に直接比例する。
}
\end{figure}


%%%%%%%%%%%%%%%%%%%%%%%%%%%%%%%%%%%%%
\subsection{オイラー離散化}

問題を離散化する第一の方法は、両方の分布$\be=\sum_{j=1}^m \b_j \de_{y_j}$と$\al_\th=\sum_{i=1}^n \a(\th)_i \de_{x_i}$が固定された位置$(x_i)_i$と$(y_j)_j$上に定義された離散分布であると仮定することである。このような位置は、観測空間全体を分割するグリッドのセル、または連続空間内の有限個の関心点の部分集合（例えば、与えられた辞書内のすべての単語に対するベクトル埋め込みの族~\citep{kusner2015word,pmlr-v51-rolet16}）を表す場合がある。この場合、パラメータ化された測度$\al_\th$は重みベクトル$\a : \th \mapsto \a(\th) \in \Si_n$を通じて完全に表現され、グリッドが大きい場合には実際には非常にスパースになりうる。この設定はいわゆるオイラー離散化法のクラスに対応する。元の形では、問題~\eqref{eq-wloss-generic}の目的関数は微分可能ではない。滑らかな最小化問題を得るために、エントロピー正則化された最適輸送を使用し、~\eqref{eq-wloss-generic}を以下で近似する
\eq{
	\umin{\th \in \Theta} \Ee_E(\th) \eqdef \MKD_\C^\epsilon(\a(\th),\b)
	\qwhereq
	\C_{i,j} \eqdef c(x_i,y_j).
}
命題~\ref{prop-convexity-dual}により、エントロピー損失関数は入力ヒストグラムに関して微分可能かつ凸であり、その勾配を持つことを想起する。

%%%%%%%%%%
\begin{proposition}[ヒストグラムに関する微分]
$\epsilon>0$に対して、
$(\a,\b) \mapsto \MKD_\C^\epsilon(\a,\b)$は凸かつ微分可能である。その勾配は
\eql{\label{eq-diff-marginals}
		\nabla \MKD_\C^\epsilon(\a,\b) = (\fD,\gD),
}
と書ける。ここで$(\fD,\gD)$は~\eqref{eq-dual-formulation}の一意な解であり、$\sum_i \fD_i=\sum_j \gD_j=0$となるように中心化されたものである。$\epsilon=0$の場合、この公式は$\MKD_\C^\epsilon$の劣微分の要素を定義し、それらが一意であれば関数は微分可能である。
\end{proposition}

$(\fD,\gD)$のゼロ平均条件は、質量保存を保証するために勾配降下法を使用する際に重要である。
%
% \todoK{
% \begin{proof}
% 	Use the dual formulation~\eqref{eq-dual-formulation}. \todoK{write me}
% \end{proof}
% }
% %%%%%%%%%%
連鎖律を用いると、$\Ee_E$は滑らかであり、その勾配は
\eql{\label{eq-diff-loss-eul}
	\nabla \Ee_E(\th) = [ \partial \a(\th) ]^\top( \fD\, ),
}
と得られる。ここで$\partial \a(\th) \in \RR^{n \times \text{dim}(\Theta)}$は写像$\a(\th)$のヤコビアン（微分）であり、$\fD \in \RR^n$はコスト行列$\C$（オイラー的設定では固定されており、特に$\th$に依存しない）に対する$\a(\th)$と$\b$の間の双対エントロピー最適輸送~\eqref{eq-dual-formulation}に付随する双対ポテンシャルベクトルである。
%
この結果を用いて、勾配降下法により$\Ee_E$を局所的に最小化できる。

%%%%%%%%%%%%%%%%%%%%%%%%%%%%%%%%%%%%%
\subsection{ラグランジュ離散化}

別のアプローチは、代わりに固定された（通常は一様な）重みを使用し、入力測度$\alpha$を点群のパラメータ化写像$x : \th \mapsto x(\th) = (x(\th)_i)_{i=1}^n \in \Xx^n$を用いた経験測度$\al_\th = \frac{1}{n}\sum_i \de_{x(\th)_i}$で近似することである。ここで$\X$はユークリッド空間であると仮定する。問題~\eqref{eq-wloss-generic}は以下のように近似される
\eql{\label{eq-lagr-sinkloss}
	\umin{\th} \Ee_L(\th) \eqdef \MKD_{\C(x(\th))}^\epsilon(\ones_n/n,\b)
	\qwhereq
	\C(x)_{i,j} \eqdef c(x(\theta)_i,y_j).
}
ここで、$\al_\th$の台が$\theta$とともに変化するため、コスト行列$\C(x(\th))$は$\th$に依存することに注意する。
%
以下の命題は、エントロピー最適輸送損失がコスト行列の滑らかな関数であることを示し、その勾配の表現を与える。

%%%%%%%%%%
\begin{proposition}[コストに関する微分]
入力ヒストグラム$(\a,\b)$を固定した場合、$\epsilon>0$に対して、写像$\C \mapsto \Rr(\C) \eqdef \MKD_\C^\epsilon(\a,\b)$は凹かつ滑らかであり、
\eql{\label{eq-diff-cost}
	\nabla \Rr(\C) = \P,
}
が成り立つ。ここで$\P$は~\eqref{eq-regularized-discr}の一意な最適解である。
%
$\epsilon=0$の場合、この公式は上勾配の集合を定義する。
\end{proposition}
%%%%%%%%%%


\todoK{
\begin{proof}
	Use the primal formulation~\eqref{eq-regularized-discr}. \todoK{write me}
	Formula~\eqref{eq-diff-positions} follows using the chain rule.
\end{proof}
}

$(\X,\Y)$が$\RR^\dim$の凸部分集合であると仮定し、~\eqref{eq-pair-discr}の形の離散測度$(\al,\be)$に対して、連鎖律を用いると$x = (x_i)_{i=1}^n  \in \X^n \mapsto \Ff(x) \eqdef  \MKD_{\C(x)}(\ones_n/n,\b)$は滑らかであり、
\eql{\label{eq-diff-positions}
	\nabla \Ff(x) = \pa{ \sum_{j=1}^m \P_{i,j} \nabla_1 \c(x_i,y_j) }_{i=1}^n \in \X^n,
}
が得られる。ここで$\nabla_1 \c$は第一変数に関する勾配である。
%
例えば、$\X=\Y=\RR^d$上で$\c(s,t) = \norm{s-t}^2$の場合、
\eql{\label{eq-diff-positions-eucl}
	\nabla \Ff(x) = 2 \pa{ \a_i x_i - \sum_{j=1}^m \P_{i,j} y_j }_{i=1}^n,
}
となる。ここで$\a_i=1/n$である。
%
定数を除けば、この勾配は$\Id-\T$であることに注意する。ここで$\T$は~\eqref{eq-baryproj}で定義された重心射影である。
%
連鎖律を用いると、ラグランジュ離散化問題~\eqref{eq-lagr-sinkloss}は滑らかであり、その勾配は
\eql{\label{eq-diff-loss-lagr}
	\nabla \Ee_L(\th) = [ \partial x(\th) ]^\top( \nabla \Ff( x(\th) ) ),
}
と得られる。ここで$\partial x(\th) \in \RR^{\text{dim}(\Theta) \times (n\dim)}$は写像$x(\th)$のヤコビアンであり、$\nabla \Ff$は$\al_\th$と$\be$の間の最適結合行列を$\P$として~\eqref{eq-diff-positions}または~\eqref{eq-diff-positions-eucl}により実装される。
%
したがって、例えば~\citep{CuturiBarycenter}で使用されているように、$\Ee_L$の局所最小解を計算するための勾配降下法を実装できる。

%%%%%%%%%%%%%%%%%%%%%%%%%%%%%%%%%%%%%
\subsection{自動微分}
\label{rem-auto-diff}

公式~\eqref{eq-diff-loss-eul}および~\eqref{eq-diff-loss-lagr}を適用する際の困難は、これらの公式が有効であるためには厳密な最適解$\fD$または$\P$を計算する必要があり、それは非常に多数のシンクホーン反復を用いてのみ許容可能な精度で達成できることである。
%
比較するヒストグラムのサイズと数が大きいような困難な状況では、単一のワッサーシュタイン距離を計算するための計算予算は通常限られており、少数のシンクホーン反復しか許容されない。そのような場合、与えられた反復で得られた値を用いて勾配~\eqref{eq-dual-formulation}を近似するよりも、逆方向自動微分を用いてシンクホーンアルゴリズムの出力を直接微分する方が通常は良い。これは~\eqref{eq-algorithmic-loss}で導入された「アルゴリズム的」シンクホーンダイバージェンスを、正則化最適輸送のエントロピーを含む~\eqref{eq-regularized-discr}の量$\MKD_\C^\varepsilon$ではなく使用し、入力（オイラーの場合はヒストグラム、ラグランジュの場合はコスト行列）を用いた単純な写像の合成として直接微分することに対応する。\S\ref{sec-regularized-cost}で導入された定義を用いると、これはそれぞれオイラーおよびラグランジュの場合において、十分大きな$L$に対して
$$\itL{\SINKHORND_{\C}}(\a(\theta),\b) \quad\text{ or }\quad \itL{\SINKHORND_{\C(x(\theta))}}(\a,\b)$$
を$\theta$に関して微分することと等価である。

シンクホーンダイバージェンスを含む汎関数の勾配を計算するコストは、汎関数自体の計算コストと同じである。このアプローチの応用については、例えば~\citep{2016-bonneel-barycoord,2017-Genevay-AutoDiff}を参照されたい。また、コスト行列に関するシンクホーン反復の微分（ラグランジュの枠組みで行われるように）に関する初期の研究として、ランキング学習への応用を含む~\citep{adams2011ranking}も参照されたい。
%
自動微分に関するさらなる詳細は~\citep{griewank2008evaluating,rall1981automatic,Neidinger10}に見出すことができ、特に勾配を計算する最も速い方法である「逆方向モード」について述べられている。
%
実装の観点から、最近のすべての深層学習Pythonフレームワークは最先端の逆方向モード微分とGPU/TPU計算のサポートを備えており~\citep{theano2016,abadi2016tensorflow,pytorch}、シンクホーン損失の大規模応用にはこれらを採用すべきである。
%
このような自動微分技術の使用を強く推奨する。なぜなら、それらは~\eqref{eq-diff-loss-eul}および~\eqref{eq-diff-loss-lagr}の計算と同じ計算量を持ち、これらの公式は主に自動微分が何を計算しているかの理論的理解を得るために有用だからである。唯一の欠点は、逆方向モード自動微分がメモリ集約的であること（メモリが反復回数に比例して増加する）である。しかし、この問題を軽減するサブサンプリング戦略が存在する~\citep{griewank1992achieving}。





%%%%%%%%%%%%%%%%%%%%%%%%%%%%%%%%%%%%%
\section{ワッサーシュタイン重心、クラスタリングおよび辞書学習}
\label{sec-bary}

教師なし学習における基本的な問題は、複数のデータ点の「平均」または「重心」を計算することである。距離空間$(\X,d)$（$d$は距離またはより一般にダイバージェンス）に属する点$(x_s)_{s=1}^S \in \X^S$の重み付き平均を定義する古典的な方法は、変分問題を解くことである
\eql{\label{eq-frechet-means}
	\umin{x \in \X} \sum_{s=1}^S \la_s d(x,x_s)^p
}
ここで$(\la_s)_s \in \simplex_S$は与えられた重みの族であり、$p$はしばしば$p=2$に設定される。
%
$\X=\RR^d$かつ$d(x,y)=\norm{x-y}_2$のとき、$p=2$では通常の線形平均$x=\sum_s \la_s x_s$の定義に、$p=1$ではより精緻なメディアン点の定義に帰着する。この形式を用いて様々な平均の概念（\emph{例：}$\X=\RR_+$上の調和平均や幾何平均）を導出できる。
%
この過程はしばしば「フレシェ」平均または「カルヒャー」平均と呼ばれる（歴史的経緯については~\citet{karcher2014riemannian}を参照）。一般的な距離$d$に対して、問題~\eqref{eq-frechet-means}は通常困難な非凸最適化問題である。幸いなことに、最適輸送距離の場合、問題は凸計画として定式化でき、存在性が証明でき、効率的な数値スキームが存在する。


%%%%%%%%%%%%
\paragraph{ワッサーシュタイン空間上のフレシェ平均。}

入力ヒストグラム$\{\b_s\}_{s=1}^S$（$b_s \in \simplex_{n_s}$）と重み$\la \in \simplex_S$が与えられたとき、ワッサーシュタイン重心は以下を最小化することで計算される
\eql{\label{eq-wass-discr}
	\umin{\a \in \simplex_n} \sum_{s=1}^S \la_s \MKD_{\C_s}(\a,\b_s),
}
ここでコスト行列$\C_s \in \RR^{n \times n_s}$を指定する必要がある。
%
典型的な設定は「オイラー的」であり、すべての重心が同じグリッド上に定義され、$n_s=n$、$\C_s=\C=\distD^p$を距離行列として設定し、以下を解く
\eq{
	\umin{\a \in \simplex_n} \sum_{s=1}^S \la_s \WassD_p^p(\a,\b_s).
}

重心問題~\eqref{eq-wass-discr}は、\citet{carlierekelandmatching}の先行するアイデアに基づき、\citet{Carlier_wasserstein_barycenter}によって任意の測度を含むより一般的な形で導入された。その提示は注~\ref{rem-bary-carlier}に委ねる。% They proved in particular uniqueness of the barycenter for $c(x,y)=\norm{x-y}^2$ over $\X=\RR^d$, if one of the input measure has a density with respect to the Lebesgue measure (and more generally under the same hypothesis as the one guaranteeing the existence of a Monge map, see Remark~\ref{rem-exist-mongemap}). ---> THIS SHOULD APPEAR IN THE GREY BOX
ヒストグラムの重心問題~\eqref{eq-wass-discr}は実際には線形計画問題である。なぜなら、各入力と重心自体との間の$S$個の結合$(\P_s)_s$を求めることができ、それらは構成上同じ行周辺を共有するよう制約されるからである。
\eq{
	\umin{\a \in \simplex_n, (\P_s \in \RR^{n \times n_s})_s } \enscond{
		\sum_{s=1}^S \la_s \dotp{\P_s}{\C_s}
	}{
		\foralls s, \P_s^\top \ones_{n_s}=\a, \P_s^\top \ones_{n} = \b_s
	}.
}
この問題は線形計画問題であるが、そのスケールにより中規模問題に対しても汎用ソルバーの使用は禁止される。双対に対する劣勾配降下法のような一次法を使用することもできる~\citep{Carlier-NumericsBarycenters}。

%%%%%%%%%%%%%%%%%%%%
\begin{rem2}{任意の測度の重心}\label{rem-bary-carlier}
	ある空間$\X$上に定義された入力測度の集合$(\be_s)_s$が与えられたとき、重心問題は以下のようになる
	\eql{\label{eq-barycenter-generic}
		\umin{\al \in \Mm_+^1(\X)} \sum_{s=1}^S \la_s \MK_{\c}(\al,\be_s).
	}
	$\X=\RR^d$かつ$c(x,y)=\norm{x-y}^2$の場合、\citet{Carlier_wasserstein_barycenter}は、入力測度の一つが密度を持てば、この重心は一意であることを示した。
	%
	問題~\eqref{eq-barycenter-generic}は、点$(x_s)_{s=1}^S \in \X^S$の重心を計算する問題を任意の測度に一般化したものと見なせる。実際、$\be_s=\de_{x_s}$が単一のディラック質量であれば、~\eqref{eq-barycenter-generic}の解は$\de_{x^\star}$であり、ここで$x^\star$は~\eqref{eq-frechet-means}を解くフレシェ平均である。
	%
	$c(x,y)=\norm{x-y}^2$の場合、重心$\al^\star$の平均は必然的に平均の重心であることに注意する。すなわち
	\eq{
		\int_\Xx x \d\al^\star(x) =  \sum_s \la_s \int_\Xx x \d\al_s(x),
	}
	であり、$\al^\star$の台は$(\al_s)_s$の台の凸包内に位置する。
	%
	入力分布を離散的なもので近似した場合（したがって~\eqref{eq-wass-discr}を代わりに解く場合）の無限次元最適化~\eqref{eq-barycenter-generic}の近似の整合性は~\citet{Carlier-NumericsBarycenters}で研究されている。
	%
	また、~\eqref{eq-barycenter-generic}を多周辺最適輸送問題として再定式化できることにも注意する。注~\ref{eq-multimarg-bary}を参照されたい。
\end{rem2}
%%%%%%%%%%%%%%%%%%%%

\begin{rem2}{ワッサーシュタイン変分問題としての$k$平均法}\label{rem-bary-kmeans}
入力測度の族$(\be_s)_s$が一つの測度$\be$のみに限られ、この測度が$\X=\RR^{\dim}$の離散有限部分集合上に台を持ち、コストが二乗ユークリッド距離である場合、重心問題
	\eql{\label{eq-barycenter-kmeans}
		\umin{\al \in \Mm_{k}^1(\X)} \MK_{\c}(\al,\be),
	}
ここで$\al$は台のサイズが高々$k$の離散測度に制約されるが、これは$\be$に対する通常の$k$平均法問題と等価であることが示せる。実際、$k$平均法問題が出力するセントロイドは解$\al$の台に対応し、その重みは各セントロイドに割り当てられた$\be$内の点の割合に対応することが容易に示せる。$\MK_{\c}$をエントロピー正則化で近似すると、混合ガウス分布のような$k$平均法のソフトクラスタリング変種に現れる平滑化された割り当てが得られることが示せる~\citep{dessein2017parameter}。
\end{rem2}

%%%%%%%%%%%%%%%%%%%%
\begin{rem2}{分布の分布と整合性}
	~\eqref{eq-barycenter-generic}を測度の（場合によっては無限の）集まりに一般化することが可能である。この問題は、確率分布の空間$\Mm_+^1(\X)$上の確率分布$M$、すなわち$M \in \Mm_+^1(\Mm_+^1(\X))$を考えることで記述される。重心は以下の解である
	\eql{\label{eq-bary-infinite}
		\umin{\al \in \Mm_+^1(\X)} \EE_M( \MK_{\c}(\al,\be) ) = \int_{\Mm_+^1(\X)} \MK_{\c}(\al,\be) \d M(\be),
	}
	ここで$\be$は$M$に従って分布する確率測度である。$M$に従って有限個$S$の入力測度$(\be_s)_{s=1}^S$を一様にランダムに抽出し、一様な重み$\la_s=1/S$に対する~\eqref{eq-barycenter-generic}の解として$\hat \be_S$を定義できる（ここで$\hat \be_S$自体がランダム測度であることに注意する）。
	%
	問題~\eqref{eq-barycenter-generic}は「離散的」な測度$M=\sum_s \la_s \de_{\be_s}$の特殊な場合に対応する。
	%
	$\MK_{\c}(\hat\be_S,\al)$のゼロへの収束（期待値でまたは高い確率で、$\al$は~\eqref{eq-bary-infinite}の一意な解）は重心の整合性に対応し、~\citep{BigotBarycenter,leGouic2016existence,bigot2012characterization}で証明されている。
	%
	これはワッサーシュタイン空間上の大数の法則として解釈できる。この結果の中心極限定理への拡張は重要な問題である。この問題の最近の定式化と特殊な場合（一次元分布とガウス測度）における解については~\citep{panaretos2016amplitude}および~\citep{agueh2017vers}を参照されたい。
\end{rem2}
%%%%%%%%%%%%%%%%%%%%


%%%%%%%%%%%%%%%%%%%%
\begin{rem2}{不動点写像}
ユークリッド空間$\X=\RR^\dim$でグラウンドコスト$\c(x,y)=\norm{x-y}^2$を扱う場合、輸送写像を用いて重心問題を研究することが可能である。実際、$\al$が密度を持つならば、注~\ref{rem-exist-mongemap}に従い、$\al$と$\al_s$の間の最適輸送写像$\T_s$を定義でき、特に$\T_{s,\sharp}\al = \al_s$が成り立つ。平均写像
\eq{
	\T^{(\al)} \eqdef \sum_{s=1}^S \la_s \T_s
}
（上記の記法は$\al$へのこの写像の依存性を明示している）は、それ自体が$\al$と$\T^{(\al)}_\sharp \al$の間の最適写像である（最適写像の正の結合はブレニエの定理（注~\ref{rem-exist-mongemap}）により凸関数の勾配の和に等しく、凸関数の和の勾配に等しいため、再びブレニエの定理により最適である）。
~\citep{Carlier_wasserstein_barycenter}で示されるように、重心問題~\eqref{eq-bary-infinite}の一次最適性条件は実際には最適測度$\al^\star$（重心）において$\T^{(\al^\star)}=\Identity_{\RR^\dim}$（恒等写像）と読め、~\citep{alvarez2016fixed}において重心$\al^\star$が不動点方程式
\eql{\label{eq-fixed-point-bary}
	G(\al) = \al
	\qwhereq
	G(\al) \eqdef \T^{(\al)}_\sharp \al,
}
の一意な解（~\citep[Theo. 2]{zemel2017fr}で明確化された正則性条件のもとで）であることが示されている。入力測度に関する穏やかな条件のもとで、\citet{alvarez2016fixed}と\citet{zemel2017fr}は$\al \mapsto G(\al)$が$\al$が重心でないとき~\eqref{eq-bary-infinite}の目的関数を厳密に減少させること、および不動点反復$\itt{\al} \eqdef G( \it{\al} )$が重心$\al^\star$に収束することを示した。
%
この不動点アルゴリズムは、最適輸送写像が閉じた形で知られている場合（\emph{例：}ガウス分布の場合）に使用できる。同じサイズの経験測度に対してこのアルゴリズムを適用すると、モンジュ写像の代わりに最適割り当ての計算が行われる。任意のサイズのより一般的な離散測度に対しても、重心射影~\eqref{eq-baryproj}を用いてスキームを適応させることができる~\citep{CuturiBarycenter}。
\end{rem2}
%%%%%%%%%%%%%%%%%%%%



%%%%%%%%%%%%
\paragraph{特殊な場合。}

一般に、~\eqref{eq-wass-discr}または~\eqref{eq-barycenter-generic}を解くことは直截的ではないが、解が陽に得られるか簡単である特殊な場合が存在する。


%%%%%%%%%%%%%%%%%%%%
\begin{rem1}{ガウス分布の重心}
	~\citep{Carlier_wasserstein_barycenter}において、二乗ユークリッドコスト$c(x,y)=\norm{x-y}^2$に対するガウス分布$\al_s = \Nn(\mean_s,\cov_s)$の重心は、それ自体がガウス分布$\Nn(\mean^\star,\cov^\star)$であることが示されている。
	%
	~\eqref{eq-dist-gauss}を用いると、重心の平均は入力の平均であることがわかる
	\eq{
		\mean^\star = \sum_s \la_s \mean_s
	}
	一方、共分散は以下を最小化する
	\eq{
		\umin{ \cov } \sum_s \la_s \Bb(\cov,\cov_s)^2,
	}
	ここで$\Bb$はビュール計量~\eqref{eq-bure-defn}である。~\citep{Carlier_wasserstein_barycenter}で研究されているように、この凸問題の一次最適性条件により$\cov^\star$は写像の一意な正定値不動点であることが示される
	\eq{
		\cov^\star = \Psi(\cov^\star)
		\qwhereq
		\Psi(\cov) \eqdef \sum_s \la_s ( \cov^{\frac{1}{2}} \cov_s \cov^{\frac{1}{2}}  )^{\frac{1}{2}},
	}
	ここで$\cov^{\frac{1}{2}}$は半正定値行列の平方根である。
	%
	この結果は~\citep{KnottSmith,RuschendorfUckelmann}から知られており、~\citep{Carlier_wasserstein_barycenter}で証明されている。
	%
	$\Psi$は厳密には縮小写像ではないが、この不動点写像を反復すること、すなわち$\itt{\cov} \eqdef \Psi( \it{\cov} )$と定義すると、実際には解$\cov^\star$に収束する。
	% has been shown by~\citet{SuvritSra} to converge to the solution $\cov^\star$.
		%
	この方法はテクスチャ合成に~\citep{2014-xia-siims}で適用されている。\citet{alvarez2016fixed}はまた、代替の写像
	\eq{
		\bar\Psi(\cov) \eqdef
		\cov^{-\frac{1}{2}}
		\Big( \sum_s \la_s ( \cov^{\frac{1}{2}} \cov_s \cov^{\frac{1}{2}}  )^{\frac{1}{2}} \Big)^2
		\cov^{-\frac{1}{2}}
	}
	の使用を提案しており、反復$\itt{\cov} \eqdef \bar\Psi( \it{\cov} )$は収束する。
	%
	これは~\eqref{eq-fixed-point-bary}で定義された不動点写像$G$がガウス分布を保存するためであり、実際に
	\eq{
		G( \Nn(\mean,\cov) ) = \Nn(\mean^\star,\bar\Psi(\cov)).
	}
	%
	図~\ref{fig-bary-gaussian}は四つの二次元ガウス分布間の重心の計算例を二つ示している。
\end{rem1}
%%%%%%%%%%%%%%%%%%%%



\begin{figure}[h!]
\centering
\begin{tabular}{@{}c@{\hspace{20mm}}c@{}}
\includegraphics[width=.3\linewidth]{barycenters-gaussians/tensor-interp-70} &
\includegraphics[width=.3\linewidth]{barycenters-gaussians/tensor-interp-98}
\end{tabular}
\caption{\label{fig-bary-gaussian}
四つの二次元ガウス分布間の重心。各ガウス分布は共分散の主軸に沿った楕円として表示され、対応する固有値に比例した伸長を持つ。
}
\end{figure}


%%%%%%%%%%%%%%%%%%%%
\begin{rem1}{一次元の場合}\label{rem-bary-1d}
	一次元分布に対して、$\Wass_p$重心は注~\ref{rem-1d-ot-generic}で詳述されているように、輸送が単調再配置であるという事実を用いてほぼ閉じた形で計算できる。
	%
	最も単純な場合は$n$点の経験測度、すなわち$\be_s = \frac{1}{n}\sum_{i=1}^n \de_{y_{s,i}}$であり、点はソート済み$y_{s,1} \leq y_{s,2} \leq \ldots$と仮定する。~\eqref{eq-1d-empirical}を用いると重心$\al_\la$も$n$点上の経験測度となる
	\eq{
		\al_\la = \frac{1}{n}\sum_{i=1}^n \de_{x_{\la,i}}
		\qwhereq
		x_{\la,i} = A_\la(x_{s,i})_s,
	}
	ここで$A_\la$は重心写像
	\eq{
	 	A_\la(x_s)_s \eqdef \uargmin{x \in \RR} \sum_{s=1}^S \la_s |x-x_{s}|^p.
	}
	である。例えば$p=2$では$x_{\la,i} = \sum_{s=1}^S \la_s x_{s,i}$となる。
	%
	一般の場合、~\eqref{eq-cumul-defn}で定義された累積関数を用いる必要があり、~\eqref{eq-wass-cumul}を用いると
	\eq{
		\foralls r \in [0,1], \quad
		\cumul{\al_\la}^{-1}(r) = A_\la( \cumul{\al_s}^{-1}(r) )_{s=1}^S,
	}
	が得られる。これは例えば、単純なソート手続きを用いて$O(n\log(n))$回の演算で$n$点未満に台を持つ離散測度間の重心を計算するために使用できる。
\end{rem1}
%%%%%%%%%%%%%%%%%%%%



%%%%%%%%%%%%%%%%%%%%
\begin{rem1}{単純な場合}
	$\T_{r,u} : x \mapsto rx+u$をスケーリングと平行移動とし、$\al_{s} = T_{r_s,u_s,\sharp} \al_0$が初期テンプレート測度をスケーリングおよび平行移動することで得られると仮定すると、重心$\al_\la$もスケーリングと平行移動により得られる\todoK{to check}
	\eq{
		\al_\la = T_{r^\star,u^\star,\sharp} \al_0
		\qwhereq
		\choice{
			r^\star = (\sum_s \la_s/r_s)^{-1},  \\
			u^\star = \sum_s \la_s u_s.
		}
	}
\end{rem1}
%%%%%%%%%%%%%%%%%%%%


%%%%%%%%%%%%%%%%%%%%
\begin{rem1}{$S=2$の場合}
	$\X=\RR^d$かつ$c(x,y)=\norm{x-y}^2$の場合（これはより一般に測地空間に拡張できる）、$S=2$つの測度$(\al_0,\al_1)$の間の重心は既に~\eqref{eq-mc-cann-interp}で導入されたマキャン補間である。$\T_\sharp \al_0=\al_1$をモンジュ写像とすると、重心$\al_\la$は$\al_\la = (\la_1 \Id + \la_2 \T)_{\sharp} \al_0$と書ける。
	%
	式~\eqref{eq-mccann-discrete}は離散の場合の計算方法を説明している。
\end{rem1}
%%%%%%%%%%%%%%%%%%%%



%%%%%%%%%%%%
\paragraph{重心のエントロピー近似。}

エントロピー平滑化を用い、ある$\epsilon>0$に対して~\eqref{eq-wass-discr}の解を以下で近似できる
\eql{\label{eq-entropic-bary}
	\umin{\a \in \simplex_n} \sum_{s=1}^S \la_s \MKD_{\C_s}^\epsilon(\a,\b_s)
}
%
これは滑らかな凸最小化問題であり、勾配降下法で解くことができる~\citep{CuturiBarycenter,GramfortPC15}。代替手段として、半双対に対する降下法（典型的には準ニュートン法）を使用することもでき~\citep{2016-Cuturi-siims}、これは重心に追加の正則化を組み込む場合、例えば与えられたノルムに関するある種の滑らかさを課す場合に有用である。
%
\citet{2015-benamou-cisc}が指摘したように、より単純かつ非常に効果的なアプローチは~\eqref{eq-entropic-bary}を（重み付き）KL射影問題として書き直すことである
\eql{\label{eq-bary-entropy-couplings}
	\umin{ (\P_s)_s } \enscond{ \sum_{s} \la_s \epsilon \KLD( \P_s|\K_s ) }{
		\foralls s, \transp{\P_s}\ones_m = \b_s, \:
		\P_1\ones_1 =  \cdots = \P_S\ones_S,
	}
}
ここで$\K_s \eqdef e^{-\C_s/\epsilon}$と表記した。ここで重心$\a$はすべての結合$\P_s \in \RR^{n \times n_s}$の行周辺に暗黙的に符号化されており、$\a = \P_1\ones_1 =  \cdots = \P_S\ones_S$である。
%
\citet{2015-benamou-cisc}で詳述されているように、シンクホーン法をこの問題に一般化でき、これは反復射影にも対応する。これは\S\ref{sec-generalized}で詳述される一般化シンクホーン法の特殊な場合とも見なせる。
%
~\eqref{eq-bary-entropy-couplings}を解く最適結合$(\P_s)_s$はスケーリング形式で以下のように計算される
\eql{\label{eq-bary-opt}
	\P_s=\diag(\uD_s)\K\diag(\vD_s),
}
そしてスケーリングは以下のように逐次的に更新される
\begin{align}\label{eq-sinkhorn-bary}
	\foralls s \in \range{1,S}, \quad \itt{\vD}_s &\eqdef \frac{\b_s}{\transp{\K}_s \it{\uD}_s}, \\
	\foralls s \in \range{1,S}, \quad  \itt{\uD}_s &\eqdef \frac{\itt{\a}}{\K_s \itt{\vD}_s}, \label{eq-sinkhorn-bary-2}\\
		\qwhereq
		\itt{\a} &\eqdef \prod_s (  \K_s \itt{\vD}_s )^{\la_s}. \label{eq-sinkhorn-bary-3}
\end{align}
これらの反復を導出する別の方法は、双対問題の変数に対して交互最小化を行うことであり、以下の命題で詳述される。

\begin{prop}
	~\eqref{eq-bary-opt}に現れる最適な$(\uD_s,\vD_s)$は$(\uD_s,\vD_s) = ( e^{\fD_s/\epsilon},e^{\gD_s/\epsilon} )$と書ける。ここで$( \fD_s,\gD_s )_s$は以下の計画の解である（その値は\eqref{eq-entropic-bary}の値と一致する）：
	\eql{\label{eq-dual-bary-entropy}
		\umax{ ( \fD_s,\gD_s )_s } \enscond{
		\sum_s \la_s \pa{
			\dotp{\gD_s}{\b_s} - \epsilon \dotp{\K_s e^{\gD_s/\epsilon}}{e^{\fD_s/\epsilon}}
		}
		}{
			\sum_s \la_s \fD_s = 0
		}.
	}
\end{prop}

\begin{proof}	~\eqref{eq-bary-entropy-couplings}にラグランジュ乗数を導入すると
	\begin{align*}
		\umin{ (\P_s)_s, \a }
		\umax{ ( \fD_s,\gD_s )_s }
		\sum_{s} \la_s \Big(
			\epsilon \KLD( \P_s|\K_s )
			+
			\dotp{\a - \P_s\ones_m}{ \fD_s } \\
		\qquad\qquad	+
			\dotp{\b_s - \transp{\P_s}\ones_m}{ \gD_s }
		\Big).
	\end{align*}
	が得られる。強双対性が成り立つため、最小と最大を交換でき、以下が得られる
		\begin{align*}
		&\umax{ ( \fD_s,\gD_s )_s }
			\sum_s \la_s \pa{
				 \dotp{\gD_s}{\b_s}
		+
		\umin{\P_s} \epsilon \KLD( \P_s|\K_s ) - \dotp{\P_s}{\fD_s \oplus \gD_s }
		} \\
	&\qquad\qquad	+ \umin{\a}
			\dotp{ \sum_{s} \la_s \fD_s }{\a}.
	\end{align*}
	$\a$に関する陽な最小化により制約$\sum_{s} \la_s \fD_s=0$が得られ、
	\begin{align*}
		\umax{ ( \fD_s,\gD_s )_s }
			\sum_s \la_s
				 \dotp{\gD_s}{\b_s}
		-
		\epsilon \KLD^*\pa{ \frac{\fD_s \oplus \gD_s}{\epsilon}|\K_s },
	\end{align*}
	となる。ここで$\KLD^*(\cdot|\K_s)$は関数$\KLD^*(\cdot|\K_s)$のルジャンドル変換~\eqref{eq-legendre}である。
	%
	このルジャンドル変換は以下のように書ける
	\eql{\label{eq-legendre-kl}
		\KLD^*(\VectMode{U}|\K) = \sum_{i,j} \K_{i,j} (e^{\VectMode{U}_{i,j}}-1),
	}
	これにより所望の公式が示される。
	%
	~\eqref{eq-legendre-kl}を示すために、この関数は分離可能であるから、以下を計算する必要がある
	\eq{
		\foralls (u,k) \in \RR_+^2, \quad
		\KLD^*(u|k) \eqdef \umax{r} ur -   \pa{ r \log(r/k) - r + k  }
	}
	その最適性条件は$u=\log(r/k)$、すなわち$r = k e^{u}$と読め、結果が得られる。
\end{proof}

すべての他の変数を固定して各$\gD_s$に関して~\eqref{eq-dual-bary-entropy}を最小化することは、~\eqref{eq-sinkhorn-bary}により閉じた形で得られる。
%
すべての$(\fD_s)_s$に関して~\eqref{eq-dual-bary-entropy}を最小化するには、~\eqref{eq-sinkhorn-bary-3}を用いて$\a$を解く必要があり、式~\eqref{eq-sinkhorn-bary-2}が導かれる。



図~\ref{fig-barycenters-images}および図~\ref{fig-barycenters-shapes}は二次元および三次元形状の補間への応用を示している。図~\ref{fig-barycenters-surfaces}は曲面上での重心の計算を示しており、グラウンドコストは測地距離の二乗である。この図では、注~\ref{rem-geod-heat}で詳述される熱の測地線近似を用いて計算が行われている。コンピュータグラフィックスおよび画像科学へのさらなる詳細と他の応用については~\citep{2015-solomon-siggraph}を参照されたい。

%
\newcommand{\FigBaryImA}[2]{\includegraphics[width=.088\linewidth]{barycenters-2d/annulus-cross-heart-2disk/shape-#1-#2}}
\newcommand{\FigBaryImLineA}[1]{\FigBaryImA{#1}{1} & \FigBaryImA{#1}{2} & \FigBaryImA{#1}{3} & \FigBaryImA{#1}{4} & \FigBaryImA{#1}{5}}
\newcommand{\FigBaryImB}[2]{\includegraphics[width=.088\linewidth]{barycenters-2d/cat-star8-thinspiral-trefle/shape-#1-#2}}
\newcommand{\FigBaryImLineB}[1]{\FigBaryImB{#1}{1} & \FigBaryImB{#1}{2} & \FigBaryImB{#1}{3} & \FigBaryImB{#1}{4} & \FigBaryImB{#1}{5}}


\begin{figure}[h!]
\centering
\fbox{
\begin{tabular}{@{}c@{}c@{}c@{}c@{}c@{}}
\FigBaryImLineA{1} \\
\FigBaryImLineA{2} \\
\FigBaryImLineA{3} \\
\FigBaryImLineA{4} \\
\FigBaryImLineA{5}
\end{tabular}
}
\fbox{
\begin{tabular}{@{}c@{}c@{}c@{}c@{}c@{}}
\FigBaryImLineB{1} \\
\FigBaryImLineB{2} \\
\FigBaryImLineB{3} \\
\FigBaryImLineB{4} \\
\FigBaryImLineB{5}
\end{tabular}
}
\caption{\label{fig-barycenters-images}
エントロピー正則化~\eqref{eq-entropic-bary}を用いた四つの入力二次元形状間の重心。バイナリ形状を表示するために、表示される画像は閾値処理された密度を示す。
%
重み$(\la_s)_s$は正方形の四隅に関して双線形である。
}
\end{figure}



% trim=left bottom right height
\newcommand{\FigBaryShapeA}[2]{\includegraphics[width=.09\linewidth,trim=130 78 110 60,clip]{barycenters-shapes/duck-spiky-moomoo_s0-double-torus/barycenter-#1-#2}}
\newcommand{\FigBaryShapeLineA}[1]{\FigBaryShapeA{#1}{0} & \FigBaryShapeA{#1}{1} & \FigBaryShapeA{#1}{2} & \FigBaryShapeA{#1}{3} & \FigBaryShapeA{#1}{4}}
%
\newcommand{\FigBaryShapeB}[2]{\includegraphics[width=.087\linewidth,trim=160 65 150 55,clip]{barycenters-shapes/mushroom-torus-hand1-trim-star/barycenter-#1-#2}}
\newcommand{\FigBaryShapeLineB}[1]{\FigBaryShapeB{#1}{0} & \FigBaryShapeB{#1}{1} & \FigBaryShapeB{#1}{2} & \FigBaryShapeB{#1}{3} & \FigBaryShapeB{#1}{4}}

\begin{figure}[h!]
\centering
%%% VOLUMETRIC SHAPES %%%
\fbox{
\begin{tabular}{@{}c@{}c@{}c@{}c@{}c@{}}
\FigBaryShapeLineA{0} \\
\FigBaryShapeLineA{1} \\
\FigBaryShapeLineA{2} \\
\FigBaryShapeLineA{3} \\
\FigBaryShapeLineA{4}
\end{tabular}
}
\fbox{
\begin{tabular}{@{}c@{}c@{}c@{}c@{}c@{}}
\FigBaryShapeLineB{0} \\
\FigBaryShapeLineB{1} \\
\FigBaryShapeLineB{2} \\
\FigBaryShapeLineB{3} \\
\FigBaryShapeLineB{4}
\end{tabular}
}
\caption{\label{fig-barycenters-shapes}
エントロピー正則化~\eqref{eq-entropic-bary}を用いた四つの入力三次元形状間の重心。重み$(\la_s)_s$は正方形の四隅に関して双線形である。
%
形状は、形状の境界内では一様で外では零であるような測度として表現される。
}
\end{figure}
ワッサーシュタイン重心の効率的な計算は、現時点でなお活発な研究テーマである~\citep{NIPS2017_6858,NIPS2018_8274}。方法論的な関心を超えて、ワッサーシュタイン重心は形状解析の分野以外でも多くの応用が見出されている。画像処理~\citep{rabin-ssvm-11}、特に色変換~\citep{2015-solomon-siggraph}（図~\ref{fig-colors}参照）；測度を要約するためのベイズ計算~\citep{srivastava2015scalable,srivastava2015wasp}；および非線形次元削減、すなわち入力測度を他の既知の測度のワッサーシュタイン重心として表現するもの~\citep{2016-bonneel-barycoord}に使用されている。これらの問題はすべて、自動微分（注~\ref{rem-auto-diff}参照）を用いて正確に最適化できる複雑な非凸目的関数をもたらす。
%
重心の計算に密接に関連する問題には、ワッサーシュタイン空間上の主成分分析の計算（例えば~\citep{SeguyCuturi,bigot2017geodesic}を参照）およびテンプレートモデルの統計的推定~\citep{boissard2015distribution}がある。重心を計算する能力は、確率測度の空間上の$k$平均法~\citep{del2016robust,ho2017multilevel}のようなより高度なクラスタリング手法を可能にする。




\begin{figure}[h!]
\centering
\begin{tabular}{@{}c@{}c@{}c@{}c@{}c@{}}
\includegraphics[width=.195\linewidth]{barycenters-surfaces/scisor-1}&
\includegraphics[width=.195\linewidth]{barycenters-surfaces/scisor-2}&
\includegraphics[width=.195\linewidth]{barycenters-surfaces/scisor-3}&
\includegraphics[width=.195\linewidth]{barycenters-surfaces/scisor-4}&
\includegraphics[width=.195\linewidth]{barycenters-surfaces/scisor-5}
\end{tabular}
\caption{\label{fig-barycenters-surfaces}
熱の測地線高速カーネル近似（注~\ref{rem-geod-heat}参照）を用いて計算された、曲面上の二つの入力測度間の重心補間。~\citep{2015-solomon-siggraph}より抜粋。
}
\end{figure}



\begin{figure}[h!]
\centering
\includegraphics[width=1\linewidth]{colors/colors-bary}
\caption{\label{fig-colors}
二つの入力画像の二つの三次元色経験ヒストグラム間の補間（ここでは簡単のため二次元色度射影のみが可視化されている）。変更されたヒストグラムは注~\ref{rem-barycenric-proj}で詳述される重心射影を用いて入力画像に適用される。~\citep{2015-solomon-siggraph}より抜粋。
}
\end{figure}



%%%%%%%%%%%%%%%%%%%%
\begin{rem1}{ワッサーシュタイン伝播}
	\citet{Solomon-ICML}で研究されているように、重心問題~\eqref{eq-wass-discr}を一般化することが可能である。ここではグラフ$\Gg$内の与えられたノード集合$U$上の分布$(\b_u)_{u \in U}$を、ノードの補集合$V$上の固定された入力分布$(\b_v)_{v \in V}$が与えられたもとで求める。未知量は、グラフのエッジを形成するすべてのノード対$(r,s) \in \Gg$間の総輸送距離を最小化することにより決定される
	\eql{\label{eq-w-propag}
		\umin{(\b_u  \in \simplex_{n_u})_{u \in U}}
			\sum_{(r,s) \in \Gg}  \MKD_{\C_{r,s}}(\b_r,\b_s),
	}
	ここでコスト行列$\C_{r,s} \in \RR^{n_r \times n_s}$はユーザーが指定する必要がある。
	%
	重心問題~\eqref{eq-wass-discr}はこの問題の特殊な場合であり、考慮されるグラフ$\Gg$は「星型」であり、$U$は他のすべての頂点$V$に接続された単一の頂点である（$\b_s$に付随する重み$\la_s$はコスト行列に吸収できる）。
	%
	各エッジ$(r,s) \in \Gg$に対して結合$\P_{r,s} \in \CouplingsD( \b_r,\b_s )$を陽に導入し、エントロピー正則化を用いると、~\eqref{eq-bary-entropy-couplings}と同様にこの問題を書き直すことができ、シンクホーン反復~\eqref{eq-sinkhorn-bary}をこの問題に拡張できる（これは\S\ref{sec-generalized}で詳述される一般化シンクホーンアルゴリズムの形にこの問題を再定式化することでも導出できる）。
	%
	\todoK{show a graphical illustration}
	%
	グラフ上のこの離散変分問題~\eqref{eq-w-propag}は、グラフを連続領域に置き換えた場合にディリクレエネルギーを定義するように一般化できる~\citep{solomon2013dirichlet}。これにより測度値調和関数の定義が導かれ、画像および曲面処理への応用が見出される。理論的解析については~\citet{Lavenant2017}、非二次的（全変動）汎関数への拡張と画像処理への応用については~\citet{vogt2017measure}も参照されたい。
\end{rem1}
%%%%%%%%%%%%%%%%%%%%



%%%%%%%%%%%%%%%%%%%%%%%%%%%%%%%%%%%%%%%%%%%%%%%%%%%%%%%%%%%%%%%%%%%%%%%%%%
%%%%%%%%%%%%%%%%%%%%%%%%%%%%%%%%%%%%%%%%%%%%%%%%%%%%%%%%%%%%%%%%%%%%%%%%%%
%%%%%%%%%%%%%%%%%%%%%%%%%%%%%%%%%%%%%%%%%%%%%%%%%%%%%%%%%%%%%%%%%%%%%%%%%%
\section{勾配流}
\label{sec-grad-flows}

滑らかな関数$\a \mapsto F(\a)$が与えられたとき、標準的な勾配降下法
\eql{\label{eq-explicit-euclidean}
	\itt{\a} \eqdef \it{\a} - \tau \nabla F(\it{\a}),
}
を使用できる。ここで$\tau$は十分小さなステップサイズである。これはいわゆる「陽的」最小化スキームに対応し、滑らかな関数$F$にのみ適用される。非滑らかな関数に対しては、代わりに「陰的」スキームを使用でき、これは近接点アルゴリズムとも呼ばれる（例えば~\citet{BauschkeCombettes11}参照）
\eql{\label{eq-implicit-euclidean}
	\itt{\a} \eqdef \Prox_{\tau F}^{\norm{\cdot}}(\it{\a}) \eqdef \uargmin{\a} \frac{1}{2}\norm{\a-\it{\a}}^2 + \tau F(\a).
}
これは~\eqref{eq-prox-eucl}で既に遭遇したユークリッド近接作用素に対応することに注意する。
%
更新~\eqref{eq-explicit-euclidean}は陽的作用素$\Id-\tau\nabla F$の反復と理解でき、一方~\eqref{eq-implicit-euclidean}は陰的作用素$(\Id+\tau\nabla F)^{-1}$を使用する。凸な$F$に対しては、反復~\eqref{eq-implicit-euclidean}は$\tau>0$のいかなる値に対しても常に収束する。

関数$F$がヒストグラムのシンプレックス$\simplex_n$上で定義されている場合、~\eqref{eq-implicit-euclidean}の$\ell^2$ノルム$\norm{\cdot}$の代わりに最適輸送距離を使用することが理にかなっており、以下を解くことになる
\eql{\label{eq-grad-flow-discr}
	\itt{\a} \eqdef \uargmin{\a} \WassD_p(\a,\it{\a})^p + \tau F(\a).
}

%%%%%%%%%%%%%%%%%%
\begin{rem2}{ワッサーシュタイン勾配流}
式~\eqref{eq-grad-flow-discr}は、反復
\eql{\label{eq-grad-flow-stepping-cont}
	\itt{\al} \eqdef \uargmin{\al} \Wass_p(\al,\it{\al})^p + \tau F(\al)
}
を定義することで任意の測度に一般化でき、ここで$F$は$\Mm_+^1(\X)$上で定義されたある関数である。この陰的時間ステッピングは連続流を構成するための有用なツールであり、形式的に極限$\tau \rightarrow 0$をとり時間$t=\tau \ell$を導入することで、$\it{\al}$は連続流$t \in \RR_+ \mapsto \al_t$を近似することを意図している。$p=2$かつ$\X=\RR^d$の特殊な場合には、形式的な計算により$\al_t$は以下の形のPDEを解くことが期待される
\eql{\label{eq-grad-flow-cont}
	\pd{\al_t}{t} = \diverg( \al_t \nabla( F'(\al_t) ) ),
}
ここで$F'(\al)$は連続関数$F'(\al) \in \Cc(\X)$であるという意味での関数$F$の微分であり、
\eq{
	F(\al+\epsilon\xi) = F(\al) + \epsilon \int_\X F'(\al) \d\xi(x) + o(\epsilon).
}
を満たす。典型的な例は$F=-\H$を用いる場合であり、ここで$\H(\al)=\KL(\al|\Ll_{\RR^d})$は$\X=\RR^\dim$上のルベーグ測度$\Ll_{\RR^d}$に関する相対エントロピーである
\eql{\label{eq-entropy-cont}
	\H(\al) = - \int_{\RR^\dim} \density{\al}(x) (\log(\density{\al}(x)) - 1) \d x
}
（$\al$が密度を持たないとき$\H(\al)=-\infty$と設定する）。このとき~\eqref{eq-grad-flow-cont}により、この負エントロピーの勾配流は線形熱拡散
\eql{\label{eq-heat}
	\pd{\al_t}{t} = \Delta \al_t,
}
であることが示される。ここで$\Delta$は空間ラプラシアンである。
%
したがって、熱拡散はディリクレエネルギー$\int_{\RR^\dim} \|\nabla \density{\al}(x)\|^2 \d x$の「古典的な」ユークリッド流（ある意味で質量振幅に関する「垂直的」移動を行う）として、あるいは代替的に最適輸送流に対するエントロピー（ある意味で質量位置に関する「水平的」移動）として解釈できる。
%
ワッサーシュタイン勾配流への関心は、Jordan、KinderlehrerおよびOttoの画期的な論文~\citep{jordan1998variational}により喚起され、これらの発展はしばしば彼らの研究にちなんで「JKO流」と呼ばれる。\citet{ambrosio2006gradient}による専門書で詳述されているように、JKO流は距離空間における勾配流の特殊な場合である。最近のサーベイ論文~\citep{santambrogio2017euclidean}も参照されたい。
%
JKO流は特に多孔質媒体方程式~\citep{otto2001geometry}、全変動流~\citep{carlier2017total}、量子ドリフト~\citep{GianazzaARMA}、または多様体上の熱発展~\citep{ErbarHeatManifold}などの非線形発展方程式の研究に使用できる。その柔軟な形式は解に対する制約を許容し、例えば密度の上限である混雑制約を\citeauthor{maury2010macroscopic}が群衆運動のモデル化に使用した~\citep{maury2010macroscopic}（レビュー論文~\citep{SantambrogioCrowdReview}も参照）。
\end{rem2}
%%%%%%%%%%%%%%%%%%%


%%%%%%%%%%%%%%%%%%
\begin{rem2}{距離空間における勾配流}
陰的ステッピング~\eqref{eq-grad-flow-stepping-cont}は、~\citep{ambrosio2006gradient}で詳述されているように、距離空間$(\Xx,\dist)$（$\dist$は距離）上で勾配流を定義するためのより一般的な形式の特殊な場合である。
%
$x \in \Xx$に対して定義されたある関数$F(x)$に対して、陰的離散最小化ステップは以下のように定義される
\eql{\label{eq-implicit-metricflow}
	\itt{x} \in \uargmin{x \in \Xx} \dist(\it{x},x)^2 + \tau F(x).
}
JKOステップ~\eqref{eq-grad-flow-stepping-cont}は確率分布の空間上でワッサーシュタイン距離を使用する場合に対応する。場合によっては、~\eqref{eq-implicit-metricflow}は$\tau \rightarrow 0$かつ$k\tau=t$のとき連続流の極限$x_t$を持つことが示せる。
%
$\Xx$がユークリッド構造も持つ場合、$F$を線形化することで陽的ステッピングが定義される
\eql{\label{eq-explicit-metricflow}
	\itt{x} = \uargmin{x \in \Xx} \dist(\it{x},x)^2 + \tau \dotp{ \nabla F(\it{x}) }{ x }.
}
陰的公式~\eqref{eq-implicit-metricflow}とは対照的に、これは通常計算が直截的であるが不安定になりうる。陰的ステップは常に安定であり、非滑らかな$F$に対しても定義されるが、通常閉じた形では得られない。
%
図~\ref{fig-gradflow-metric}は関数$F(x) = \norm{x}^2$を$\Xx=\RR^2$上で、距離$d(x,y)=\norm{x-y}_p=(|x_1-y_1|^p+|x_2-y_2|^p)^{\frac{1}{p}}$に対して$p$のいくつかの値で示したこの概念を図解している。
%
陽的スキーム~\eqref{eq-explicit-metricflow}は$p=1$と$p=+\infty$で不安定であり、$p=1$では軸に沿ったステップ（座標ごとの降下）を与える。
%
対照的に、陰的スキーム~\eqref{eq-implicit-metricflow}は安定である。特に$p=1$の場合、二つの座標が等しいとき、次のステップが対角方向に作用することに注意されたい。
\end{rem2}
%%%%%%%%%%%%%%%%%%


\begin{figure}[h!]
\centering
\begin{tabular}{cc}
\includegraphics[width=.45\linewidth]{gradflow-metric/explicit}&
\includegraphics[width=.45\linewidth]{gradflow-metric/implicit}\\
Explicit & Implicit
\end{tabular}
\caption{\label{fig-gradflow-metric}
距離$d(x,y)=\norm{x-y}_p$に対する$p$のいくつかの値での、$\Xx=\RR^2$上の関数$f(x) = \norm{x}^2$を最小化するための陽的および陰的勾配流の比較。
}
\end{figure}


%%
\begin{rem1}{粒子系を用いたラグランジュ離散化}\label{rem-lagr-discretization}
~\eqref{eq-grad-flow-discr}の有限次元問題は、測度の空間上の流れ~\eqref{eq-grad-flow-stepping-cont}のオイラー離散化と解釈できる。
%
問題を離散化する別の方法として、いわゆる粒子系を用いたラグランジュ法では、解を代わりに時間とともに移動する（離散的な）経験測度としてパラメータ化し、その測度の位置（重みではなく）が関心のある変数となる。実際には、時間でインデックスされた粒子の動的な点群$\al_t = \frac{1}{n}\sum_{i=1}^n \de_{x_i(t)}$を考えることができる。
%
元の問題~\eqref{eq-grad-flow-discr}は、点の動力学を規定する$n$個の連立常微分方程式$X(t)=(x_i(t))_i \in \X^n$に置き換えられる。
%
エネルギー$F$が離散測度に対して有限であれば、単に$\Ff(X)=F( \frac{1}{n}\sum_{i=1}^n \de_{x_i} )$と定義できる。典型的な例は線形関数$F(\al) = \int_\X V(x) \d\al(x)$と二次相互作用$F(\al) = \int_{\X^2} W(x,y) \d\al(x)\d\al(y)$であり、この場合それぞれ
\eq{
	\Ff(X) = \frac{1}{n} \sum_i V(x_i) \qandq \Ff(X) = \frac{1}{n^2} \sum_{i,j} W(x_i,x_j).
}
を使用できる。
%
密度を持つ測度に対してのみ有限である一般化エントロピーのような関数に対しては、点群を密度に変換するための密度推定量を適用する必要があり、これにより$n \rightarrow +\infty$のとき$F$と整合的な関数$\Ff(x)$も定義できる。典型的な例は~\eqref{eq-entropy-cont}で定義されたエントロピー$F(\mu) = \H(\al)$であり、整合的な推定量（定数項を除いて）は最近傍までの距離の対数の和として得られる
\eql{\label{eq-entropy-discretized}
	\Ff(X) = \frac{1}{n}\sum_{i} \log( d_X(x_i) )
	\qwhereq d_X(x) = \umin{x' \in X, x' \neq x} \norm{x-x'};
}
ノンパラメトリックエントロピー推定量のレビューについては~\citet{beirlant1997nonparametric}を参照されたい。
%
十分小さなステップサイズ$\tau$に対して、$\X=\RR^\dim$と仮定すると、ワッサーシュタイン距離$\Wass_2$は点上のユークリッド距離と一致する。すなわち$|t-t'|$が十分小さければ$\Ww_2(\al_t,\al_{t'})=\norm{X(t)-X(t')}$である。したがって勾配流は位置に関するユークリッド流$X'(t) = -\nabla \Ff(X(t))$と等価であり、これは時刻$t_k=\tau k$に対して~\eqref{eq-explicit-euclidean}と同様に陽的オイラーステップで離散化される
\eq{
	\itt{X} \eqdef \it{X} - \tau \nabla \Ff(\it{X}).
}
%
図~\ref{fig-flow-lagr}は線形関数とエントロピー汎関数の和に対するこのような離散化された陽的発展の例を示しており、フォッカー・プランク方程式の離散化版が得られる。
%
この特殊な線形フォッカー・プランク方程式の場合、確率的PDEの方法に頼ることも可能であり、ガウスドリフトを持つ単一のランダム粒子を発展させることで数値的に近似できることに注意する。これらのスキーム（いわゆるランジュバン・モンテカルロ法）の定常分布への収束は、ワッサーシュタイン距離の観点から定量化できる。例えば~\citep{dalalyan2017user}を参照されたい。\todoK{Detailed this much more.}
%
関数$\Ff$が滑らかでない場合、~\eqref{eq-implicit-euclidean}と同様に陰的オイラーステップを用いて離散化すべきである。すなわち
\eq{
	\itt{X} \eqdef \Prox_{\tau \Ff}^{\norm{\cdot}}(\it{X}) \eqdef \uargmin{Z \in \X^n} \frac{1}{2}\norm{Z-\it{X}}^2 + \tau \Ff(Z).
}
を考える。
%
線形関数$F(\al) = \int_X V(x) \d\al(x)$の最も単純な場合、流れは各粒子$x_i(t)$に対して独立に作用し、関数$V$に対する通常のユークリッド流$x_i'(t)=-\nabla V(x_i(t))$に対応する（そして流れの積分曲線に沿った密度の移流PDEである）。
\end{rem1}

\newcommand{\FigLagrFlow}[1]{\includegraphics[width=.19\linewidth]{lagrangian-flows/lagrange-flow-#1}}


\begin{figure}[h!]
\centering
\begin{tabular}{@{}c@{}c@{}c@{}c@{}c@{}}
\FigLagrFlow{1} &
\FigLagrFlow{3} &
\FigLagrFlow{5} &
\FigLagrFlow{7} &
\FigLagrFlow{10} \\
$t=0$ & $t=0.2$  & $t=0.4$ & $t=0.6$ & $t=0.8$
\end{tabular}
\caption{\label{fig-flow-lagr}
関数$F(\al) = \int V \d\al - H(\al)$（$V(x)=\norm{x}^2$）に対するラグランジュ離散化を用いた勾配流発展の例。
%
エントロピーは~\eqref{eq-entropy-discretized}を用いて離散化されている。
%
極限定常分布はガウス分布である。
}
\end{figure}


%%%%%%%%%%%%%%%%%%
\begin{rem2}{測地的凸性}
勾配流に関連する重要な概念は、ワッサーシュタイン-2幾何学に関する汎関数$F$の凸性、すなわちワッサーシュタイン測地線（注~\ref{rem-displacement}で示される変位補間）に沿った$F$の凸性である。そのような関数に対するワッサーシュタイン勾配流（連続時間）は存在し、一意であり、$\tau \rightarrow 0$のときの離散ステッピング~\eqref{eq-grad-flow-stepping-cont}の極限である。これは$t \rightarrow +\infty$のとき固定された定常分布に収束する。
%
エントロピーは測地的凸関数の典型的な例であり、線形関数$F(\al) = \int_\X V(x) \d\al(x)$および凸関数$V : \X \rightarrow \RR$、$W : \X \times \X \rightarrow \RR$に対する二次相互作用関数$F(\al) = \int_{\X \times \X} W(x,y) \d\al(x) \d\al(y)$もそうである。
%
線形関数は古典的な意味で凸であるが、二次相互作用関数はそうでない場合があることに注意する。
%
典型的な例は$W(x,y)=\norm{x-y}^2$であり、これは負半正定核（定義~\ref{def-negativedefinitekernel}参照）であるため、$F(\al)$は通常の意味で凹な関数に対応する（一方で測地的には凸である）。
%
\citet{mccann1997convexity}の重要な結果として、$\X=\RR^\dim$上の形式$F(\al) = \int_{\RR^\dim} \phi(\density{\al}(x)) \d x$の一般化「エントロピー」関数は、$\phi$が凸であり、$\phi(0)=0$、$t \rightarrow +\infty$のとき$\phi(t)/t \rightarrow +\infty$、かつ$s \mapsto s^{\dim} \phi(s^{-\dim})$が凸減少であるならば、測地的に凸である。
\end{rem2}
%%%%%%%%%%%%%%%%%%%

結果として得られる離散化された流れの数値解法に関する重要な文献があり、ここではいくつかの代表的な出版物のみを挙げる。一次元問題に対しては、最適輸送が逆累積関数において二次汎関数であるため（注~\ref{rem-1d-ot-generic}参照）、非常に精密なソルバーが開発されている：\citet{kinderlehrer1999approximation,blanchet2008convergence,agueh2013one,Matthes1D,blanchet2012optimal}。高次元では、有限要素法および有限体積法を用いて取り組むことができる：\citet{CarrilloFiniteVolume,burger2010mixed}。代替的なソルバーはラグランジュスキーム（すなわち粒子系）を用いて得られる：\citet{carrillo2009numerical,JDB-JKO,Westdickenberg2010}。
%
もう一つの方向は、連続的な対応物の一部の性質を維持する離散流（典型的には離散グリッドまたはグラフ上）を探すことである。~\citet{MielkeCVPDE,ErbarDCDS,ChowHuangLiZhou2012,Maas2011}を参照されたい。


エントロピー正則化に基づくオイラー離散化問題~\eqref{eq-explicit-euclidean}を解くための近似的アプローチは~\citet{2015-Peyre-siims}で最初に提案され、~\citet{2016-chizat-sinkhorn}で改良され、~\citet{2017-carlier-SIMA}で理論的に解析された。
%
エントロピー正則化を用いると、問題~\eqref{eq-grad-flow-discr}は$G = \iota_{\it{\a}}$と設定し$F$を$\tau F$に置き換えたとき~\eqref{eq-generalized-ot-regul}の形になる。したがって、~\citet{2015-Peyre-siims}で最初に提案されたように、反復~\eqref{eq-gen-sinkh}を用いて$\itt{\a}$を近似できる。$\epsilon \rightarrow 0$のときのこのスキームの収束は~\citet{2017-carlier-SIMA}で証明されている。
%
図~\ref{fig-jko}はこの方法で計算された発展の例を示している。
%
勾配流の機械学習への興味深い応用は、密度のある動力学モデルを最もよくモデル化する基礎関数$F$を学習することである。この学習はエントロピー正則化された輸送と自動微分（注~\ref{rem-auto-diff}参照）を用いた滑らかな非凸最適化を解くことで達成できる。~\citet{hashimoto2016learning}を参照されたい。

時間と空間の両方で離散化された勾配流の収束を解析することは一般に困難である。距離を定義する線形計画問題の多面体的性質により、ステップサイズを小さくしすぎると「ロッキング」現象（分布が動かなくなり発展しなくなるため、ステップサイズは小さすぎてはならない、~\citep{maury201713}で議論されている）が生じる。一次元における勾配流の離散化法の収束解析については~\citep{Matthes1D,matthes2017convergent}を参照されたい。


\begin{figure}[h!]
\centering
\begin{tabular}{@{}c@{\hspace{1mm}}c@{\hspace{1mm}}c@{\hspace{1mm}}c@{\hspace{1mm}}c@{}}
\includegraphics[width=.19\linewidth]{jko/congestion/holes-potential}&
\includegraphics[width=.19\linewidth]{jko/congestion/holes-kappa10-1}&
\includegraphics[width=.19\linewidth]{jko/congestion/holes-kappa10-2}&
% \includegraphics[width=.19\linewidth]{jko/congestion/holes-kappa10-5}&
\includegraphics[width=.19\linewidth]{jko/congestion/holes-kappa10-10}&
\includegraphics[width=.19\linewidth]{jko/congestion/holes-kappa10-20}\\
\includegraphics[width=.19\linewidth]{jko/surfaces/cosw}&
\includegraphics[width=.19\linewidth]{jko/surfaces/evol-1}&
\includegraphics[width=.19\linewidth]{jko/surfaces/evol-2}&
\includegraphics[width=.19\linewidth]{jko/surfaces/evol-3}&
\includegraphics[width=.19\linewidth]{jko/surfaces/evol-4}\\
$\cos(w)$ & $t=0$ & $t=5$ & $t=10$ & $t=20$
\end{tabular}
\caption{\label{fig-jko}
ドリフト$V$と混雑項を伴う勾配流発展の例（~\citet{2015-Peyre-siims}より）。$F(\al) = \int_\X V(x) \d \al(x) + \iota_{\leq \kappa}(\density{\al})$である。
}
\end{figure}



\S\ref{sec-unbalanced}で詳述される非均衡最適輸送距離に対しても勾配流を計算することが可能である。これにより質量の生成または消滅を許容する発展が得られ、多くの物理的、生物学的、または化学的現象をモデル化する上で不可欠である。% Figure~\ref{fig-ot-grad-flows} shows
非均衡勾配流の例は、細胞成長のためのヘレ・ショー模型~\citep{PerthameTumor}であり、~\citep{gallouet2017jko,marino2017tumor}で理論的に研究されている。このような非均衡勾配流も一般化シンクホーンアルゴリズム~\citep{2016-chizat-sinkhorn}を用いて近似できる。


%\begin{figure}[h!]
%\centering
%\includegraphics[width=.8\linewidth]{jko/heleshaw_profile}
%\caption{\label{fig-ot-grad-flows}
%Unbalanced OT gradient flow to solve the Hele-Shaw equation (from~\citep{2016-chizat-sinkhorn}). \todoK{detail the functional involved}
%}
%\end{figure}


%%%%%%%%%%%%%%%%%%%%%%%%%%%%%%%%%%%%%
%%%%%%%%%%%%%%%%%%%%%%%%%%%%%%%%%%%%%
%%%%%%%%%%%%%%%%%%%%%%%%%%%%%%%%%%%%%
%%%%%%%%%%%%%%%%%%%%%%%%%%%%%%%%%%%%%
\section{最小カントロヴィッチ推定量}
\label{sec-mke}

ある未知の分布からの離散サンプル$(x_i)_{i=1}^n \subset \Xx$が与えられたとき、目的はパラメトリックモデル$\th \mapsto \al_\th \in \Mm(\Xx)$を観測された経験入力測度$\be$にフィッティングすることである
\eql{\label{eq-density-fitting}
	\umin{\th \in \Theta} \Ll(\al_\th,\be)
	\qwhereq
	\be = \frac{1}{n} \sum_i \de_{x_i},
}
ここで$\Ll$は離散分布と「連続」（任意の）分布の間のある「損失」関数である（図~\ref{fig-density-fitting}参照）。

$\al_\th$がルベーグ測度（または他の固定された参照測度）に関する密度$\density{\th} \eqdef \density{\al_\th}$を持つ場合、最尤推定量（MLE）は以下を解くことで得られる
\eq{
	\umin{\th} \Ll_{\text{MLE}}(\al_\th,\be) \eqdef -\sum_i \log(\density{\th}(x_i)).
}
これは、$x_i$がある$\bar\be$のi.i.d.サンプルであると仮定すると、カルバック・ライブラー損失の経験的対応物を使用することに対応する。なぜなら
\eq{
	\Ll_{\text{MLE}}(\al,\be) \overset{n \rightarrow +\infty}{\longrightarrow} \KL(\al|\bar\be).
}
であるからである。

\begin{figure}[h!]
\centering
\includegraphics[width=.4\linewidth]{density-fitting/schematic-fitting}
\caption{\label{fig-density-fitting}
密度フィッティング問題~\ref{eq-density-fitting}の模式図。
}
\end{figure}


\newcommand{\fPF}{h}

このMLEアプローチは多くの場合に最適な推定手続きをもたらすことが知られている（例えば~\citet{owen2001empirical}参照）。しかし、特異な分布を推定する場合、典型的には$\al_\th$が密度を持たない場合（$\Ll_{\text{MLE}}(\al_\th,\be) = +\infty$となる）、または$(x_i)_i$がある特異な$\bar\be$からのサンプルである場合（$\KL(\al_\th|\bar\be)$が有限であるためには$\al_\th$が$\be$と同じ台を共有すべきであるが、この台は通常未知である）にはうまく機能しない。もう一つの問題は、実用上の関心が高い場合のいくつかにおいて、密度$\density{\th}$が計算できない（または計算が困難すぎる）ことである。

両方の問題（特異性と未知の密度）が生じる典型的な設定は、いわゆる生成モデルであり、パラメトリック測度は固定された参照測度$\zeta \in \Mm(\Zz)$の押し出しとして書かれる
\eq{
	\al_\th = \fPF_{\th,\sharp} \zeta \qwhereq \fPF_\th : \Zz \rightarrow \Xx,
}
ここで押し出し作用素は定義~\ref{defn-pushfwd}で導入されている。空間$\Zz$は通常低次元であるため、$\al_\th$の台は低次元の「多様体」に沿って局在化し、結果として得られる密度は高度に特異的である（ルベーグ測度に関する密度を持たない）。
%
さらに、この密度の計算は通常計算困難であるが、$\al_\th$からのi.i.d.サンプルの生成は$x_i=\fPF_\th(z_i)$を計算することで達成される。ここで$(z_i)_i$は$\zeta$からのi.i.d.サンプルである。

このような困難なシナリオに対処するためには、MLE汎関数$\Ll_{\text{MLE}}$の代わりに弱い距離を使用する必要があり、それは双対形式として
\eql{\label{eq-dual-loss}
	\Ll(\al,\be) \eqdef
	\umax{(\f,\g) \in \Cc(\X)^2}
	\enscond{ \int_{\X} \f(x) \d\al(x) + \int_{\X} \g(x) \d\be(x) }{ (\f,\g) \in \Potentials}.
}
と書く必要がある。\S\ref{sec-dual-norms}で示される双対ノルムは
\eq{
	\Potentials = \enscond{(\f,-\f)}{\f \in B},
}
を課すことに対応し、最適輸送~\eqref{eq-dual-generic}は$\Potentials = \Potentials(\c)$を~\eqref{eq-dfn-pot-dual}で定義されたように設定する。

固定された$\th$に対して、このような損失関数を用いて~\eqref{eq-density-fitting}で最小化すべきエネルギーを評価することは、半離散最適輸送を解くことに対応し、これは第~\ref{c-algo-semidiscr}章の焦点である。$\th$に関してエネルギーを最小化することははるかに複雑であり、典型的には高度に非凸である。

$\Ee(\th) = \Ll(\al_\th,\be)$を評価する際の~\eqref{eq-dual-loss}の解を$\f_\th$と表記すると、劣勾配は公式
\eql{\label{eq-grad-wass-loss-dual}
	\nabla \Ee(\th) = \int_\Xx [\partial \fPF_\th(x)]^\top \nabla \f_\th(x) \d\al_\th(x),
}
を用いて得られる。ここで$\partial \fPF_\th(x) \in \RR^{\text{dim}(\Theta) \times \dim}$は$\th \in \RR^{\text{dim}(\Theta)} \mapsto \fPF_\th(x)$の（$\th$に関する）微分であり、$\nabla \f_\th(x)$は$\f_\th$の（$x$に関する）勾配である。
%
この公式は数値的に使用するのが困難である。まず、半離散問題の解である連続関数$\f_\th$を計算する必要があるからである。\S\ref{sec-entropy-ot-mmd}で示されるように、最適輸送損失に対しては確率的最適化を用いて達成できるが、高次元では困難である。もう一つの選択肢は、このポテンシャルに対してパラメトリックな形式、例えばRKHSにおける展開（\citet{genevay2016stochastic}）や深層ネットワーク近似（\citep{WassersteinGAN}）を課すことである。しかし、これは理論的にはまだ解析されていない重要な近似誤差をもたらす。
%
最後の問題は、双対ポテンシャル$\f_\th$の勾配$\nabla \f_\th$の計算を必要とするため、数値的に不安定であることである。

最適輸送損失に対しては、以下の等価な問題に対する主最適結合を計算する方が、代替的な勾配公式が得られる：
\eql{\label{eq-ot-couplings-fitting}
	\MK_\c(\al_\th,\be) = \umin{\ga \in \Mm(\Zz \times \Xx)}
		\enscond{
			\int_{\Zz \times \Xx} c(\fPF_\th(z),x) \d\ga(z,x)
		}{
			\gamma \in \Couplings(\zeta,\be)
		}.
}
ここで考えている半離散の場合、最小化すべき目的関数は実際には以下のように分解できることに注意する
\eql{\label{eq-ot-couplings-fitting-bis}
	\umin{(\ga_i)_{i=1}^n} \sum_{i=1}^n \int_\Zz c(\fPF_\th(z),x_i) \d \ga_i(z)
	\qwhereq
	\sum_{i=1}^n \ga_i = \zeta, \quad
	\int_\Zz \d\ga_i(z)  = \frac{1}{n},
}
ここで各$\ga_i \in \Mm_+^1(\Zz)$である。
%
~\eqref{eq-ot-couplings-fitting-bis}を解く最適な$(\gamma_{\th,i})_i$が得られれば、$\Ee(\th)$の勾配は以下のように計算される
\eq{
	\nabla \Ee(\th) = \sum_{i=1}^n \int_\Zz [\partial \fPF_\th(z)]^\top \nabla_1 c(\fPF_\th(z),x_i) \d \ga_i(z),
}
ここで$\nabla_1 c(x,y) \in \RR^\dim$は$x \mapsto c(x,y)$の勾配である。~\eqref{eq-grad-wass-loss-dual}とは対照的に、この公式は双対最適輸送問題の解であるポテンシャルの勾配の計算を含まないことに注意する。

$\Ll=\MK_\c$を用いて得られるこの推定量のクラスは、しばしば「最小カントロヴィッチ推定量」と呼ばれ、~\citep{bassetti2006minimum}で最初に導入された。~\citep{CanasRosasco}も参照されたい。これは制限ボルツマンマシンの訓練のために~\citep{CuturiBoltzman}で生成モデルの文脈で使用され、近似ベイズ計算と組み合わせて~\citep{bernton2017inference}で使用されている。深層ネットワークを用いたこれらの計算の近似は、GAN~\citep{WassersteinGAN}とVAE~\citep{Bousquet2017}の両方の深層生成モデルの訓練に使用されている。~\citep{2017-Genevay-AutoDiff,2017-Genevay-gan-vae,salimans2018improving}も参照されたい。パラメトリックモデルフィッティングに対するシンクホーンダイバージェンスの使用は、形状マッチングおよびレジストレーションに日常的に使用されている。~\citep{gold1998new,chui2000new,myronenko2010point,2017-feydy-miccai}を参照されたい。


\begin{rem1}{距離学習と転移学習}
	機械学習への応用において、最適輸送関連手法の成功はデータの幾何学を捉えた適切なコスト$c(x,y)$の選択に大きく依存することを強調しておく。
	%
	多くの種類のデータをユークリッド空間に埋め込むことが可能であるが（例えば単語埋め込みについては~\citep{mikolov2013efficient}参照）、多くの場合、距離の何らかの適応または最適化が必要である。
	%
	教師あり課題に対する距離学習は古典的な問題であり（例えば~\citep{MAL-019,weinberger2009distance}参照）、学習パイプラインで最適輸送距離が使用される場合のグラウンド距離$c(x,y)$の学習に拡張されている~\citep{CuturiGroundMetric2014}（~\citealt{ZenICPR14,WangECCV12OLD,huang2016supervised}も参照されたい）。
	%
	また、最適結合$\P$の観測からコスト行列を学習するという関連する逆問題も挙げておく。これは低ランク事前分布を用いて正則化できる~\citep{dupuy2016estimating}。
	%
	関連する問題には転移学習~\citep{pan2010survey}とドメイン適応~\citep{glorot2011domain}があり、ここでは訓練済みの機械学習パイプラインを新しいデータセットに適応させたい。この問題は最適輸送技術を用いてモデル化および解くことができる。~\citep{courty2017optimal,courty2017joint}を参照されたい。
\end{rem1}
